% Nota de trabalho (não sai no PDF). Todos os números vêm do banco, com n = 46
% (7 itens vazados excluídos via not_leaked) e preços do Together AI (24/10/2026).
% Reprodução:
%   uv run python -m src.evaluation.run_validity
%   uv run python -m src.evaluation.run_fidelity
%   uv run --with transformers python -m src.evaluation.run_custo
%   uv run --with transformers python -m src.evaluation.run_tcr [--all-arms]
%   uv run --with transformers python -m src.evaluation.export_figuras

\chapter{Resultados}\label{chap:resultados}

\section{Visão geral}\label{sec:res-visao-geral}

Este capítulo segue a ordem das medidas da Seção~\ref{sec:medidas}. Primeiro vêm a validade (Seção~\ref{sec:res-validade}), o custo (Seção~\ref{sec:res-custo}) e a compressão (Seção~\ref{sec:res-compressao}) que sustentam as conclusões principais. Depois é mostrada a fidelidade, que é exploratória (Seção~\ref{sec:res-fidelidade}), e uma análise do modelo especializado, que reúne os braços do modelo pequeno e os braços exploratórios (Seção~\ref{sec:res-modelo-especializado}). A Tabela~\ref{tab:res-panorama} apresenta de forma resumida os braços principais nas três medidas que os comparam diretamente.

\begin{table}[htb]
	\Caption{\label{tab:res-panorama} Panorama dos braços}%
	\IBGEtab{}{%
		\small
		\begin{tabular}{lllrrrr}
			\toprule
			Braço & Modelo & Saída & Validade & Itens cobertos & Fidelidade & Custo \\
			\midrule \midrule
			A1  & DeepSeek~V4~Pro & XML & 97,8\% & 46 & 3,77 & 5.561 \\
			A1g & GLM~5.2 & XML & 93,5\% & 46 & 3,32 & 6.386 \\
			A2  & DeepSeek~V4~Pro & DSL & 62,3\% & 34 & 3,67 & 4.311 \\
			A2g & GLM~5.2 & DSL & 52,9\% & 35 & 3,49 & 5.505 \\
			A3  & Qwen~7B & DSL & 58,7\% & 27 & --- & 874 \\
			A3m & Qwen~7B & DSL & 0,0\% & 0 & --- & --- \\
			A4  & Qwen~7B$^{\dagger}$ & DSL & 87,0\% & 40 & 2,78 & 591 \\
			\midrule
			A2r & DeepSeek~V4~Pro & DSL & 87,0\% & 43 & 3,85 & 3.241 \\
			\bottomrule
		\end{tabular}%
	}{%
	\Fonte{Elaborado pelo autor}%
	\Nota{Validade: proporção das 138 gerações (46 itens, 3 repetições) aceitas pelo esquema XSD. Itens cobertos: itens com ao menos uma saída válida, de 46. Fidelidade: média das notas do juiz de 0 a 5 nos itens cobertos, medida exploratória (Seção~\ref{subsec:fidelidade}); o juiz não foi aplicado ao A3 e ao A3m. Custo: US\$ por milhão de saídas válidas, sem \textit{cache}; para o A4, no cenário com acréscimo de 50\% pelo adaptador. $^{\dagger}$~com o adaptador treinado. O A2r é exploratório.}%
}
\end{table}

A geração direta de XML teve a maior validade: 97,8\% no A1 e 93,5\% no A1g. Com os mesmos modelos, pedir a DSL no lugar do XML reduziu a validade para 62,3\% e 52,9\%, o contrário do que o protocolo inicial previa. A maioria dessas falhas estava nas referências a elementos já escritos, e o A2r, que acrescenta a referência por rótulo, elevou a validade do mesmo modelo para 87,0\%. Mesmo assim, ele continua abaixo da geração direta (Seção~\ref{sec:res-validade}).

O modelo pequeno especializado, o A4, teve a mesma validade do A2r, 87,0\%, sem a gramática no \textit{prompt}. A igualdade, porém, é uma coincidência: dos itens em que algum dos dois falhou, só um é comum aos dois braços, de modo que eles não esbarram em uma mesma limitação da notação (Seção~\ref{sec:res-validade}). Ela também não permite afirmar que os dois são equivalentes, porque com 46 itens a incerteza da diferença é maior que a margem fixada (Seção~\ref{sec:protocolo}). O custo por saída válida do A4 é de US\$~591 por milhão, contra US\$~3.241 do A2r e US\$~5.561 do A1, mesmo no cenário de preço mais desfavorável ao A4 (Seção~\ref{sec:res-custo}).

Na fidelidade, que é exploratória, o A2r e o A1 ficaram no mesmo nível, e o A4 ficou cerca de um ponto abaixo de ambos, uma diferença maior que a margem fixada (Seção~\ref{sec:res-fidelidade}). O resultado do A4 é, portanto, uma troca: a mesma validade do método no modelo de grande porte, a um custo várias vezes menor, com menor fidelidade.

\section{Validade}\label{sec:res-validade}

\subsection{Validade por braço}

A Figura~\ref{fig:validade} mostra a validade de cada braço, isto é, a proporção das 138 gerações que produziram um XML aceito pelo esquema. Os braços se dividem em três grupos. A geração direta de XML ficou acima de 90\%, com 97,8\% no A1 e 93,5\% no A1g. O A4 e o A2r ficaram em 87,0\%. Os braços que pedem a DSL pela gramática no \textit{prompt}, sem a referência por rótulo, ficaram entre 52,9\% e 62,3\%: o A2, o A2g e o A3.

% Validade sintática por braço. Barras ordenadas por magnitude — a comparação é
% de ordenação, não de série temporal, e ordem alfabética esconderia o achado.
% `bar shift=0pt` em cada série: sem isso o pgfplots desloca lateralmente cada
% \addplot e as barras deixam de coincidir com o rótulo do eixo.
% Legenda horizontal acima da área de plotagem, para não cobrir barra alguma.
\begin{figure}[htb]
	\Caption{\label{fig:validade} Validade sintática por braço experimental}
	\centering
	\begin{tikzpicture}
		\begin{axis}[
			width=0.86\textwidth, height=6.4cm,
			ybar, bar width=17pt, bar shift=0pt,
			ymin=0, ymax=108,
			ylabel={Saídas válidas no XSD (\%)},
			symbolic x coords={A1,A1g,A4,A2r,A2,A3,A2g,A3m},
			xtick={A1,A1g,A4,A2r,A2,A3,A2g,A3m},
			nodes near coords, nodes near coords style={font=\scriptsize, text=black, /pgf/number format/fixed, /pgf/number format/precision=1},
			ymajorgrids, grid style={draw=black!12},
			axis lines*=left, tick align=outside,
			tick label style={font=\small}, label style={font=\small},
			legend style={font=\small, at={(0.5,1.04)}, anchor=south, draw=none, legend columns=3, column sep=1.2em},
			legend cell align=left,
		]
			\addplot+[draw=none, fill=corXML]
				table[x=braco, y=xsd, col sep=comma] {figuras/dados/validade_xml.csv};
			\addlegendentry{XML direto}
			\addplot+[draw=none, fill=corSFT]
				table[x=braco, y=xsd, col sep=comma] {figuras/dados/validade_sft.csv};
			\addlegendentry{DSL, 7B com ajuste fino}
			\addplot+[draw=none, fill=corDSL]
				table[x=braco, y=xsd, col sep=comma] {figuras/dados/validade_dsl.csv};
			\addlegendentry{DSL por \textit{prompting}}
		\end{axis}
	\end{tikzpicture}
	\Fonte{Elaborado pelo autor}
	\Nota{Percentual sobre as 159 gerações de cada braço. Nos braços de DSL a
	validade XSD coincide com a taxa de \textit{parsing}, pois a transpilação é
	determinística e sempre produz XML válido a partir de DSL aceita.}
\end{figure}


O A3m não produziu nenhuma saída válida, e o motivo não é a dificuldade da notação. Sem o adaptador e sem a gramática no \textit{prompt}, o modelo não tem como saber o que é a DSL, e escreveu BPMN XML nas 138 gerações, o formato que conhece do pré-treinamento. Em 81 delas, o XML passou do limite de 2.048 \textit{tokens}, dimensionado para a DSL, e foi truncado. O zero do A3m mostra que o \textit{prompt} mínimo, sozinho, não transmite a notação; é com esse mesmo \textit{prompt} que o A4 chega a 87,0\%, e a diferença entre os dois é apenas o adaptador.

Entre os três braços do modelo pequeno, a validade passa de 0\% sem a notação, a 58,7\% com a gramática no \textit{prompt} (A3) e a 87,0\% com a notação aprendida no ajuste (A4). O A4 ainda usa um \textit{prompt} 598 \textit{tokens} menor que o do A3, porque dispensa a gramática (Seção~\ref{sec:bracos}).

\subsection{Comparações com margem de equivalência}

A Tabela~\ref{tab:res-validade-contrastes} mostra os dois contrastes de validade definidos na Seção~\ref{sec:protocolo}, com a margem de $\pm$10 pontos percentuais.

\begin{table}[htb]
	\Caption{\label{tab:res-validade-contrastes} Contrastes de validade}%
	\IBGEtab{}{%
		\begin{tabular}{lrrl}
			\toprule
			Contraste & Diferença & IC 95\% & Veredito \\
			\midrule \midrule
			A4 $-$ A2r & 0,0 & [$-$13,8; 13,0] & inconclusivo \\
			A2r $-$ A1 & $-$10,9 & [$-$21,0; $-$1,4] & inconclusivo \\
			\bottomrule
		\end{tabular}%
	}{%
	\Fonte{Elaborado pelo autor}%
	\Nota{Diferença média da validade por item, em pontos percentuais, pareada nos 46 itens. Intervalo de confiança por \textit{bootstrap} com 10.000 reamostragens. Margem de equivalência de $\pm$10 pontos, fixada antes do cálculo dos intervalos.}%
}
\end{table}

Os dois contrastes são inconclusivos, mas por motivos diferentes. Entre o A4 e o A2r não há diferença observada, mas o intervalo vai de $-$13,8 a 13,0 pontos e ultrapassa a margem dos dois lados, de modo que não é possível afirmar que os dois são equivalentes. Entre o A2r e o A1, o intervalo inteiro está abaixo de zero: o A2r é menos válido que a geração direta de XML. A perda pode, porém, ser menor que 10 pontos, porque o intervalo começa em $-$21,0 e termina em $-$1,4, e por isso também não é possível afirmar que ela ultrapassa a margem.

\subsection{Causas de invalidade}

A Tabela~\ref{tab:res-falhas} classifica as saídas inválidas dos braços de DSL pela causa da rejeição.

\begin{table}[htb]
	\Caption{\label{tab:res-falhas} Causas de invalidade nos braços de DSL}%
	\IBGEtab{}{%
		\begin{tabular}{lrrrrr}
			\toprule
			Braço & Inválidas & Referência & Outra sintaxe & Bloco aberto & Truncada \\
			\midrule \midrule
			A2  & 52 & 38 & 14 & 0 & 0 \\
			A2g & 65 & 41 & 23 & 1 & 0 \\
			A2r & 18 & 8 & 10 & 0 & 0 \\
			A3  & 57 & 0 & 57 & 0 & 0 \\
			A3m & 138 & 0 & 57 & 0 & 81 \\
			A4  & 18 & 12 & 0 & 6 & 0 \\
			\bottomrule
		\end{tabular}%
	}{%
	\Fonte{Elaborado pelo autor}%
	\Nota{Contagem sobre as 138 gerações de cada braço. Referência: referência a um identificador nunca declarado, ou marca de referência (\texttt{\#}) em posição que a gramática não aceita. Bloco aberto: a saída termina com um bloco sem fechamento. No A3m, todas as saídas são XML, e não DSL.}%
}
\end{table}

No A2 e no A2g, a referência é a principal causa de falha: 38 das 52 saídas inválidas do A2 e 41 das 65 do A2g. Em parte delas, o modelo referencia um identificador que nunca declarou. Nas demais, a marca de referência aparece onde a gramática não a aceita, em geral numa tentativa de declarar o identificador depois de já tê-lo usado. Nos dois casos, o problema é o mesmo: para voltar a um elemento, o modelo precisava tê-lo marcado com um identificador quando o escreveu pela primeira vez, antes de saber que voltaria a ele. Foi esse diagnóstico que motivou o A2r (Seção~\ref{sec:bracos}), e com a referência por rótulo as falhas de referência do mesmo modelo caíram de 38 para 8. As 8 que restaram são referências a identificadores nunca declarados, em 4 itens.

O A4 e o A2r têm 18 saídas inválidas cada, mas falham de formas diferentes. O A4 falha em 6 itens, sempre nas três repetições, o que é esperado da decodificação gulosa. Suas falhas são 12 referências a identificadores nunca declarados e 6 blocos não fechados. A referência continua sendo o principal problema do A4, o que é coerente com o treino, que só contém a referência por identificador (Seção~\ref{sec:transpilacao}). O A2r falha em 8 itens, a maioria em apenas uma ou duas das três repetições. Só um item, o pmo\_32, falha nos dois braços. Se os dois esbarrassem em uma construção que a DSL não consegue expressar, as falhas cairiam nos mesmos itens.

O A3 não tem nenhuma falha de referência. Suas 57 saídas inválidas são erros de sintaxe de outra natureza, como palavras-chave inventadas e construções fora da gramática, o que indica que o modelo pequeno sem ajuste não domina a notação nem com a gramática no \textit{prompt}.

\section{Custo}\label{sec:res-custo}
\section{Compressão}\label{sec:res-compressao}
\section{Fidelidade (exploratória)}\label{sec:res-fidelidade}
\section{O modelo especializado}\label{sec:res-modelo-especializado}
\section{Síntese}\label{sec:res-sintese}
